\section{Twelve positive unknowns: a tighter separating norm}
\label{sec:positive}

The construction combines the separating-root encoding of package~23 with
rational Pell saturation from packages~19--22. Saturation removes the extra
factor $H$ in package~23's strong multiplier. A sharper separating weight
then deliberately cancels leading terms in four conjugate factors. The
result improves package~20's degree-$1145$, $123$-operation equation
\cite{repo19,repo20,repo21,repo23}. All statements in this section use the
entire fixed coefficient tuple specified elsewhere in this report.

\subsection{An explicit polynomial}

The twelve supplied unknowns are
\[
 J,F,\alpha,z,h,i,T,s,w,\eta,\zeta,t\in\Z_{>0}.
\]
All other symbols are polynomial abbreviations. The fixed ports are $\beta,K_0,M_C,M_F$, with $\beta=B-1$ and $M_F$ the shifted high
mask. Define
\begin{align*}
 q&=\beta J+1,& C&=q-F-\alpha,& X&=wq,& Y&=sq^3,& E&=XY,\\
 k&=\eta+\zeta,&c&=kY+\eta,&a&=E+Y,&H&=4a+3,&\Delta&=a^2+H,\\
 \ell&=EkY,& A&=1+\ell(\ell+k),&G&=1+\Delta c^2,&b_0&=X+ac.
\end{align*}
The transport and index expressions are
\begin{align*}
 R&=((q-1)(q-F)+\alpha)(q^2-1)+(M_C+qM_F)J,\\
 N_k&=k-hE-R,\\
 N_t&=(K_0+w)C+q-F-z(q-1).
\end{align*}
The normalized strong and auxiliary expressions are
\begin{align*}
 I&=ic^2,& S&=\Delta I,&Q&=S^2,&D&=1+IS,&d&=Q-1,\\
 v&=cT,& b&=c+RD,&e&=Q(Dv^2+b^2)-1,&B_1&=2Qvb.
\end{align*}
In particular $\Delta=(a+2)^2-1$ and $D=1+\Delta I^2$.
Choose one of the following three weights:
\begin{align}
 M_{1081}&=1+8I(T+2REic),\label{eq:m1081}\\
 M_{1073}&=I\bigl(2T+R[E+3(X+Y)]ic\bigr),\label{eq:m1073}\\
 M_{1069}&=I\bigl(2T+R[E+2(X+Y)+sq^2]ic\bigr).
 \label{eq:m1069}
\end{align}
Set
\[
 m=HM,\qquad N=1+2m,\qquad W=Ht+Nb_0.
\]
For formal roots $\tau^2=A$, $g^2=G$, put
\begin{align}
 \mathcal F(\tau,g)
 &=\bigl[d(W-m\tau-Ng)^2-H^2e\bigr]^2-H^4B_1^2D,
 \label{eq:formalF}\\
 \Omega&=\prod_{\tau=\pm\sqrt A,\,g=\pm\sqrt G}
                   \mathcal F(\tau,g).\label{eq:omega}
\end{align}
The product is an integer polynomial, without additional unknowns.
The requested single equation is
\begin{equation}
 \boxed{\quad P=N_kN_t(1+\Omega)-1=0.\quad}
 \label{eq:positiveP}
\end{equation}
The corresponding three-equation system is
\begin{equation}
 N_k=1,\qquad N_t=1,\qquad \Omega=0.
 \label{eq:positiveSystem}
\end{equation}

For clarity, no algebraic-number or resultant operation is needed to
compute $\Omega$. The following scalar circuit uses only integer arithmetic:
\begin{align*}
 r_0&=W^2+N^2G,&\Theta&=4W^2N^2G,&a_0&=m^2A,\\
 C_0&=d(r_0+a_0)-H^2e,&k_f&=4d^2a_0,&K_c&=k_fr_0,\\
 Z_0&=H^4B_1^2D,&x_0&=C_0^2+\Theta d^2-K_c-Z_0,
   &x_1&=2C_0d-k_f,\\
 n_0&=x_0^2+\Theta x_1^2-4Z_0K_c,
   &n_1&=2x_0x_1-4Z_0k_f,\\
 \Omega&=n_0^2-\Theta n_1^2.
\end{align*}
To verify this identity, write $\theta=2WN g$, so $\theta^2=\Theta$.
After eliminating $\tau$, the conjugate factor is
$(C^2-K-Z_0)^2-4KZ_0$, where
$C=C_0+d\theta$ and $K=K_c+k_f\theta$.
Its even and odd parts are $n_0,n_1$, respectively.
The system compares $n_0^2=\Theta n_1^2$ directly, saving the final
subtraction. Shared intermediate values and the factorization of $M$ are
charged exactly as in the attached JSON circuits.

\begin{theorem}\label{thm:positive12}
For each weight in \eqref{eq:m1081}--\eqref{eq:m1069}, the positive
integer solvability of \eqref{eq:positiveP} or
\eqref{eq:positiveSystem} is equivalent to blank halting of the calibrated
K215 machine. Consequently ZFC proves that absence of a positive integer
solution is equivalent to $\operatorname{Con}(\mathrm{ZFC})$.
The exact sizes are as follows; every binary addition, subtraction, or
multiplication, including one by a fixed numeral, is charged.
\end{theorem}

\begin{center}
\begin{tabular}{lrrrr}
\hline
Weight & System degree & System operations & Equation degree & Operations\\
\hline
$M_{1081}$ &1072&115&1081&120\\
$M_{1073}$ &1064&116&1073&121\\
$M_{1069}$ &1060&118&1069&123\\
\hline
\end{tabular}
\end{center}

The sum-of-squares equation
$(N_k-1)^2+(N_t-1)^2+\Omega^2=0$ has degrees
$2144,2128,2120$ and operation counts $123,124,126$, respectively.
The theorem concerns ordinary mathematics formalizable in ZFC. It does
not assert a new kernel-checked end-to-end reduction or prove consistency.

\subsection{The integrality lemma that removes the extra multiplier}

\begin{lemma}[Rational Pell saturation]\label{lem:saturation}
Let $s_0\ge2$ and $V$ be integers. If $y\in\mathbb Q$ and
\[
 s_0^2V^2-(s_0^2-1)y^2=1,
\]
then $y\in\Z$.
\end{lemma}
\begin{proof}
Write $y=u/v_0$ in lowest terms with $v_0>0$. Its equation implies
$v_0^2\mid s_0^2-1$. The integer
$D_0=(s_0^2-1)/v_0^2$ is positive and nonsquare, since otherwise
$s_0^2-1$ would be a square. The pairs
$(s_0,v_0)$ and $(s_0|V|,|u|)$ solve $x^2-D_0z^2=1$ with positive
ordinates; $u=0$ or $V=0$ is impossible. Write them as
$(\chi_j,\psi_j)$ and $(\chi_l,\psi_l)$ in the powers of the fundamental
positive integer-coordinate Pell unit, where $j,l\ge1$.

We have $\chi_j\mid\chi_l$. The addition identities imply
\[
 \chi_{r+2j}\equiv-\chi_r\pmod{\chi_j},\qquad
 \chi_{2j-r}\equiv-\chi_r\pmod{\chi_j}.
\]
Reduce $l$ modulo $2j$, then reflect to $0\le r\le j$.
Strict growth gives $0<\chi_r<\chi_j$ for $r<j$, so the divisible
case requires $r=j$. Thus $l$ is an odd multiple of $j$.
The Pell multiplication formula gives $\psi_j\mid\psi_l$, hence
$v_0\mid u$. Their coprimality forces $v_0=1$.
\end{proof}

This is the special feature of the auxiliary norm. Rationality of a root
of a general nonmonic quadratic would not imply integrality. Packages
19--22 establish equivalent saturation lemmas
\cite{repo19,repo20,repo21}; applying the lemma before recovering the main
congruence removes package~23's extra $H$ without changing any canonical
compiler witness.

\subsection{Separation bounds before recovering roots}

First, $\Theta$ is divisible by $4$, so $\Omega$ is a square modulo $4$.
At a zero of \eqref{eq:positiveP}, its three integer factors are units.
The last cannot be $-1$: that would require $\Omega=-2$, which is not a
square modulo $4$. Therefore
\[
 \Omega=0,\qquad N_k=N_t=\epsilon\in\{1,-1\}.
\]
The system has $\epsilon=1$. The inherited weak transport bootstrap is
valid for both signs and gives, before any eliminated root is recovered,
\begin{equation}
 \begin{gathered}
 C\ge0,\quad q\ge16,\quad X\ge16,\quad Y\ge q^3\ge4096,\\
 0<R<q^4-2,\qquad E\ge q^4>R+2.
 \end{gathered}\label{eq:weakbounds}
\end{equation}
These are coefficient-dependent statements from the authentic blank
interface, not assertions for arbitrary substituted ports
\cite{normalized,repo20,repo21}.

A vanished factor in \eqref{eq:omega} supplies coherent roots satisfying
\begin{align}
 \tau^2&=A,&g^2&=G,&f^2&=D,\nonumber\\
 d y^2&=e-B_1f=Q(vf-b)^2-1,\label{eq:rootrelations}\\
 t&=y+M\tau+N\gamma,&\gamma&=(g-b_0)/H.
 \label{eq:linearroot}
\end{align}
The roots here are initially algebraic. In particular, no integrality
or sign of $\gamma$ has been assumed.

For the first weight, all complex auxiliary branches satisfy
\[
 |y|<2\bigl(v(a+2)I+b\bigr).
\]
Indeed $\sqrt D<(a+2)I$ and $Q\ge2$ give this bound directly from
\eqref{eq:rootrelations}. Since
$a+2<2E$, $c<2kY$, $\sqrt A>\ell>Ec/2$, and $D<4E^2I^2$,
\[
 \frac{|y|}{\sqrt A}
 <8TI+\frac4E+16REi^2c^3<M_{1081}.
\]

Here is the sharper estimate used for the other weights.
Since $\Delta>R+2$ and $I^2\ge c^4$,
\[
 Q-b=\Delta I^2(\Delta-R)-c-R>0.
\]
Thus $b\le d$; also $\sqrt D<S<d$.
Put $Z=v\sqrt D+b>1$. For every auxiliary conjugate,
\[
 |y|^2\le\frac{QZ^2+1}{d},\qquad
 \left[Z\left(1+\frac1d\right)\right]^2-
       \frac{QZ^2+1}{d}=\frac{QZ^2-d}{d^2}>0.
\]
Consequently
\begin{equation}
 |y|<Z+v+1<cT\bigl((a+2)I+1\bigr)+c+RD+1.
 \label{eq:sharpy}
\end{equation}
Let $\rho=(1+1/Y)/E$. We have $c/\ell<\rho$, and
$(a+3)\rho<2$, because
$(a+3)/E<9/8$ and $1+1/Y<17/16$.
Furthermore
\begin{align*}
 (a+2)^2/E&=E+2Y+4+Y/X+4/X+4/E<E+3Y+5,\\
 (a+2)^2\rho&<E+3Y+X+9<E+3(X+Y)-2.
\end{align*}
For $C_1=E+3(X+Y)$, dividing \eqref{eq:sharpy} by $\sqrt A$
therefore gives
\[
 \frac{|y|}{\sqrt A}
 <2TI+R(C_1-2)i^2c^3+2
 \le 2TI+RC_1i^2c^3=M_{1073}.
\]
The last inequality uses $Ri^2c^3\ge1$.

For $C_2=E+2(X+Y)+sq^2$, use the exact identity
\[
 (a+2)^2\rho=E+2Y+X+6+Y/X+5/X+4/Y+8/E+4/(EY).
\]
Since $Y/X=sq^2/w\le sq^2$ and
$5/X+4/Y+8/E+4/(EY)<1$, the right side is less than
$E+2Y+X+sq^2+7<C_2-2$. The identical argument proves
$|y|<M_{1069}\sqrt A$.

Finally $G>A$: $c>kY$ and $\Delta-E^2>2EY$ imply
\[
 \Delta c^2>k^2Y^2\Delta>
 k^2(E^2Y^2+EY)=A-1.
\]
For each of the three weights, the two uniform separation inequalities are
\begin{equation}
 |y|<M\sqrt A,\qquad
 \frac{N\sqrt G}{H}>2M\sqrt G>2M\sqrt A.
 \label{eq:separate}
\end{equation}

\subsection{Rationality, integrality, and the compiler interface}

Work in a finite normal splitting field over $\mathbb Q$ containing all
roots in \eqref{eq:rootrelations}. If an automorphism changed the sign of
$g$, it would change $N\gamma$ by magnitude $2N\sqrt G/H$.
The change in $M\tau$ has magnitude at most $2M\sqrt A$, and the change
in $y$ has magnitude strictly less than $2M\sqrt A$, using the uniform
bound for both conjugates. Their sum cannot cancel the former change by
\eqref{eq:separate}, contradicting the fixed rational value $t$ in
\eqref{eq:linearroot}. Hence $g$ is rational. With $g$ fixed, a change in
$\tau$ would similarly require $2M\sqrt A<2M\sqrt A$.
Thus $\tau$ is rational. Equation \eqref{eq:linearroot} makes $y$
rational, and $B_1>0$ makes $f=(e-dy^2)/B_1$ rational.

Since $g^2=G$, $\tau^2=A$, and $f^2=D$ are integers, $g,\tau,f$ are
integers. The integer $V=vf-b$ and $S=\Delta I\ge2$ now satisfy the
hypotheses of Lemma~\ref{lem:saturation}, so $y$ is an integer.
Both $H\gamma=g-b_0$ and $N\gamma=t-M\tau-y$ are integers, and
$N=1+2HM$ gives $\gcd(H,N)=1$. Bezout's identity proves
$\gamma\in\Z$.

The recovered equations are precisely the normalized first, main, strong,
and auxiliary Pell equations. The inherited signed-coordinate theorem
permits signed $\gamma,\tau,f,y,T$, and conversion to its ordinary strong
coordinate is $i_{\rm ord}=\Delta i$. That theorem restores the unit sign
$\epsilon=1$ and extracts a blank K215 halting computation
\cite{normalized,repo20,repo21}. In particular, the sign of the main root
is not assumed in the reconstruction.

Conversely, take a canonical positive normalized blank compiler witness.
For each chosen weight, set
\[
 t=y+M\tau+N\gamma>0.
\]
One factor of \eqref{eq:omega} vanishes and $N_k=N_t=1$.
No additional divisibility restriction on $i$ and no adjustment of the
auxiliary Pell period is required. This proves the existential equivalence
in Theorem~\ref{thm:positive12}; composing with the calibrated K215
refutation equivalence and negating existential solvability gives its
claimed ZFC theorem.

\subsection{Uniform exact degrees and the deliberate cancellation}

Scale the twelve supplied variables by a formal parameter $\lambda$ and
write a subscript $0$ for the leading homogeneous part. The shared degrees
are
\[
 \begin{array}{c|rrrrrrrrrrr}
 \text{expression}&E,H&Y&c&R&\Delta&I&Q,d&D&b&A,G&M,N\\
 \hline
 \text{degree}&6&4&5&4&12&11&46&34&38&22&27,33.
 \end{array}
\]
Put
\[
 \mathcal B=R_0E_0^2i^2c_0^4,\qquad
 \mathcal D=R_0E_0Y_0i^2c_0^4.
\]
These have degrees $38$ and $36$.
The two formal main-quotient roots solve
\[
 H^2\gamma^2+2Hb_0\gamma+(b_0^2-G)=0.
\]
The coefficient degrees are $12,17,16$, with
$(b_0^2-G)_{16}=-H_0c_0^2\ne0$. Thus the two roots have orders $5$ and
$-1$; the larger has leading term $-2\ell_0/H_0$.
Also each auxiliary root has expansion
\[
 y=\varepsilon(b-vf)+O(\lambda^{-8}),\qquad \varepsilon\in\{1,-1\},
\]
since $Q/(Q-1)=1+O(\lambda^{-46})$, while $vf$ has order $23$.
These expansions can be taken in a Laurent-series extension of the
rational function field; all leading radicands are squares there.

For $M_{1081}$, $M\sqrt A$ has leading term $16\mathcal B$ and
$N\gamma$ on its large branch has leading term $-64\mathcal B$.
The exact factorization is
\begin{equation}
 \Omega=d^8H^{16}
  \prod\bigl(t-y-M\tau-N\gamma\bigr),
 \label{eq:16factor}
\end{equation}
with sixteen coherent branches, counting multiplicity. Their leading
magnitudes are $15,17$ four times each, and $47,49,79,81$ twice each,
in units of $\mathcal B$. No leading term vanishes. Hence
\[
 \deg\Omega=8\cdot46+16\cdot6+16\cdot38=1072.
\]
The leading coefficient multiplier is
$15^4 17^4(47\cdot49\cdot79\cdot81)^2$, which is nonzero.

For $M_{1073}$, the positive formal first root satisfies
$\sqrt A-\ell=O(\lambda)$, and
\[
 M\sqrt A-b=Ri^2c^3(C_1\ell-\Delta c)+O(\lambda^{28}).
\]
Here the exact polynomial identity
\[
 C_1\ell-\Delta c
 =(EY+3EX-Y^2-4E-4Y-3)c-E(E+3X+3Y)\eta
\]
has unique leading form $E_0Y_0c_0$ of degree $15$.
Consequently $M\sqrt A-b$ has leading form $\mathcal D$ of degree $36$:
its terms of degrees $38$ and $37$ cancel identically.
Now $N\gamma$ on the large branch has leading term $-4\mathcal B$;
on the small branch it has order $32$.
In \eqref{eq:16factor}, four factors have degree $36$ and leading
magnitude $\mathcal D$, four have degree $38$ and magnitude
$2\mathcal B$, and the remaining eight have leading multiples
$2,4,4,6$ of $\mathcal B$, each twice. The signed product is positive:
\begin{align*}
 \Omega_{1064}
 &=589824\,d_0^8H_0^{16}\mathcal B^{12}\mathcal D^4,\\
 1064&=8\cdot46+16\cdot6+12\cdot38+4\cdot36.
\end{align*}
Writing $q_0=\beta J$, $k_0=\eta+\zeta$ gives the fully substituted form
\begin{equation}
 \Omega_{1064}=9\cdot2^{48}i^{48}k_0^{96}w^{76}s^{176}
                     q_0^{652}(q_0-F)^{16}.\label{eq:top1073}
\end{equation}
It is nonzero for every authentic slice $\beta>0$.
The product's leading form multiplies this by
\[
 [-hw s q_0^4]\,[w(q_0-F-\alpha)-zq_0].
\]
For example, the coefficient of
$J^{673}hz i^{48}\eta^{96}w^{77}s^{177}$ is
$9\cdot2^{48}\beta^{673}>0$. This proves exact degree $1073$ uniformly,
not merely for diagnostic numerical specializations.

Finally, for $M_{1069}$ replace $C_1$ by $C_2$ above. Then
\[
 C_2\ell-\Delta c
 =(2EX+Esq^2-Y^2-4E-4Y-3)c-EC_2\eta.
\]
Its unique leading term is $E_0(sq_0^2)c_0$, of degree $14$.
Thus the four cancellation factors have degree $35$, rather than $36$.
All other factors retain their previous degrees and leading multipliers.
The norm has exact degree $1060$, and its leading form is
\[
 \Omega_{1060}=9\cdot2^{48}i^{48}k_0^{96}w^{76}s^{176}
                     q_0^{648}(q_0-F)^{16}.
\]
The corresponding product coefficient at
$J^{669}hz i^{48}\eta^{96}w^{77}s^{177}$ is
$9\cdot2^{48}\beta^{669}>0$, proving exact degree $1069$.
The degrees of $N_k,N_t$ are $7,2$ in every case. The sum-of-squares
claims follow because $\Omega^2$ is its unique highest-degree summand.

The independent finite checks compare the literal circuits against
conjugate-product evaluators, verify all gates and ports are live, and
compute full univariate specializations over two finite fields to detect
the cancellations. They corroborate the formulas and operation counts;
the preceding uniform arguments establish the mathematical equivalence
and exact degrees. They do not construct an accepting K215 computation.
